\documentclass[a4paper]{article}

\usepackage{graphicx}
\usepackage{amsmath,amssymb}
\usepackage{minted}
\usepackage{multirow}
\usepackage{float}
\usepackage{todonotes}
\usepackage{forest}
\usepackage{xstring}
\usepackage{xcolor}
\usepackage[most]{tcolorbox}
\usepackage{xparse}
\usepackage{natbib}

\usetikzlibrary{graphs,graphdrawing}
\usegdlibrary{layered}

% for external graphviz
\input{/home/donkarlo/Dropbox/repo/nd_ai_project/src/nd_ai/knoweldge_representation/language/formal/computer/programming/tex/latex/code_clips/external_graphviz.tex}

% for tags
\input{/home/donkarlo/Dropbox/repo/nd_ai_project/src/nd_ai/knoweldge_representation/language/formal/computer/programming/tex/latex/parts/control_sequence/kinds/semtag/semtag.tex}

% for links
\input{/home/donkarlo/Dropbox/repo/nd_ai_project/src/nd_ai/knoweldge_representation/language/formal/computer/programming/tex/latex/code_clips/seehere.tex}

\title{}
\date{}

\begin{document}
    \maketitle
    The efficient coding hypothesis was proposed by Horace Barlow in 1961 as a theoretical model of sensory neuroscience in the brain. Within the brain, neurons communicate with one another by sending electrical impulses referred to as action potentials or spikes.
    
    Barlow hypothesized that the spikes in the sensory system formed a neural code for efficiently representing sensory information. By efficient it is understood that the code minimized the number of spikes needed to transmit a given signal. This is somewhat analogous to transmitting information across the internet, where different file formats can be used to transmit a given image. Different file formats require different numbers of bits for representing the same image at a given distortion level, and some are better suited for representing certain classes of images than others. According to this model, the brain is thought to use a code which is suited for representing visual and audio information which is representative of an organism's natural environment .
    \seeherefooter{https://en.wikipedia.org/wiki/Efficient_coding_hypothesis}
    \bibliographystyle{plainnat}
    \bibliography{/home/donkarlo/Dropbox/repo/nd_shared_research/bibliography/bibliography.bib}
\end{document}